\documentclass[11pt,a4paper]{article}

\usepackage[UTF8]{ctex}
\usepackage{amsmath,amssymb,amsthm}
\usepackage{bm}
\usepackage{booktabs}
\usepackage{graphicx}
\usepackage{geometry}
\usepackage{hyperref}
\usepackage{natbib}
\usepackage{xcolor}

\geometry{margin=2.5cm}
\hypersetup{colorlinks=true,linkcolor=blue,citecolor=blue,urlcolor=blue}
\graphicspath{{../word2vec/figures/}}

\newcommand{\todo}[1]{\textcolor{red}{【待补：#1】}}
\newcommand{\figfile}[4]{%
  \begin{figure}[htbp]
    \centering
    \IfFileExists{../word2vec/figures/#2}{%
      \includegraphics[width=#1]{#2}%
    }{%
      \fbox{\parbox[c][4cm][c]{0.85\linewidth}{\centering\color{red}%
      图尚未生成：\texttt{\detokenize{#2}}（训练后用 visualize.py 生成）}}%
    }%
    \caption{#3}\label{#4}
  \end{figure}%
}

\title{实验四：词向量\\\large word2vec 的梯度推导·中文词向量的训练、评测与可视化}
\author{姓名：李贵龙 \quad 学号：202411581243}
\date{\today}

\begin{document}
\maketitle

\begin{abstract}
本次实验分为理论与编程两部分。理论部分补全 Lecture 3 第 35 页略去的推导：
在朴素 softmax 下求损失对上下文向量 $v_k$ 的偏导，理解梯度“期望减观测”的含义，
并估算朴素 softmax 每一步的计算代价。编程部分在中文维基百科语料上完成
“预处理与分词 $\to$ 训练 skip-gram + 负采样词向量 $\to$ 中文词相似度与最近邻评测
$\to$ 超参数对比 $\to$ PCA/t-SNE 可视化”的完整流程。
训练使用官方 C 代码 \texttt{tmikolov/word2vec}；skip-gram 与负采样分别由
\citet{mikolov2013efficient} 与 \citet{mikolov2013distributed} 提出。
\end{abstract}

\tableofcontents
\newpage

% ============================================================
\section*{零、先读这一节}
\addcontentsline{toc}{section}{零、先读这一节}
% ============================================================

\textbf{记号约定}（与讲义一致，与 CS224n 相反）：$u_i$ 是词 $i$ 作为中心词时的向量，
$v_i$ 是词 $i$ 作为上下文词时的向量，$u_i,v_i\in\mathbb{R}^d$。
对一个训练对（中心词 $t$，上下文词 $c$），朴素 softmax 下的条件概率与损失为
\begin{equation}
P(k\mid t)=\frac{\exp(u_t\cdot v_k)}{\sum_{k'\in V}\exp(u_t\cdot v_{k'})},
\qquad
y=-\log P(c\mid t).
\label{eq:loss}
\end{equation}
把所有上下文向量按行堆成矩阵 $V\in\mathbb{R}^{|V|\times d}$（第 $k$ 行是 $v_k^\top$）；
记 $\hat y\in\mathbb{R}^{|V|}$ 为预测分布（第 $k$ 个分量是 $P(k\mid t)$），
$e_c$ 为第 $c$ 个分量为 $1$ 的 one-hot 向量。词向量的基础概念可参考
\citet{jurafsky2024slp} 第 6 章。

% ============================================================
\section{第一部分 理论（40 分）}
% ============================================================

\subsection{题 1(a)：对中心词向量求导（10 分）}

先展开对数，把损失写成两项之和。记配分函数
$Z=\sum_{k\in V}\exp(u_t\cdot v_k)$，则
\begin{equation}
y=-\log\frac{\exp(u_t\cdot v_c)}{Z}
=-\bigl(u_t\cdot v_c-\log Z\bigr)
=-u_t\cdot v_c+\log\sum_{k\in V}\exp(u_t\cdot v_k).
\label{eq:y-expand}
\end{equation}

对 $u_t$ 逐项求导。第一项 $-u_t\cdot v_c=-\sum_{i=1}^d u_{t,i}v_{c,i}$，
由 $\partial(u_t\cdot v_c)/\partial u_t=v_c$，得
\begin{equation}
\frac{\partial(-u_t\cdot v_c)}{\partial u_t}=-v_c .
\end{equation}

第二项 $\log Z$ 用链式法则：外层函数是 $\log(\cdot)$，内层函数是
$Z(u_t)=\sum_{k}\exp(u_t\cdot v_k)$。因为
\begin{equation}
\frac{\partial Z}{\partial u_t}
=\sum_{k\in V}\frac{\partial\exp(u_t\cdot v_k)}{\partial u_t}
=\sum_{k\in V}\exp(u_t\cdot v_k)\,v_k,
\end{equation}
所以
\begin{equation}
\frac{\partial\log Z}{\partial u_t}
=\frac{1}{Z}\frac{\partial Z}{\partial u_t}
=\sum_{k\in V}\frac{\exp(u_t\cdot v_k)}{Z}v_k
=\sum_{k\in V}P(k\mid t)\,v_k .
\end{equation}

合起来即得
\begin{equation}
\boxed{\;\frac{\partial y}{\partial u_t}
=-v_c+\sum_{k\in V}P(k\mid t)\,v_k\;}
\label{eq:grad-u}
\end{equation}

\textbf{矩阵形式。}
由 $V$ 的行是 $v_k^\top$ 可知
$\sum_{k}P(k\mid t)v_k=V^\top\hat y$，而 $v_c=V^\top e_c$，因此
\begin{equation}
\frac{\partial y}{\partial u_t}
=V^\top\bigl(\hat y-e_c\bigr)
\in\mathbb{R}^{d}.
\label{eq:grad-u-matrix}
\end{equation}

\subsection{题 1(b)：对上下文向量求导（18 分）}

由式 \eqref{eq:y-expand}，$y$ 对 $v_k$ 的依赖来自两部分：
第一项 $-u_t\cdot v_c$ 与对数配分项 $\log Z$。

\paragraph{(i) 第一项 $-u_t\cdot v_c$ 的贡献。}
$-u_t\cdot v_c$ 中只出现 $v_c$，与其它 $v_k\,(k\neq c)$ 无关，因此
\begin{equation}
\frac{\partial(-u_t\cdot v_c)}{\partial v_k}
=\begin{cases}
-u_t, & k=c,\\[2pt]
0, & k\neq c .
\end{cases}
\end{equation}
这正是“第一项 $-\,u_t\cdot v_c$ 只在 $k=c$ 时有贡献”的原因。

\paragraph{(ii) 对 $\log Z$ 使用链式法则。}
把每个 $v_k$ 都看作自变量：
\begin{itemize}
  \item \textbf{外层函数}：$f(g)=\log g$，其导数为 $f'(g)=1/g$，其中 $g=Z$；
  \item \textbf{内层函数}：$g(v_k)=Z=\sum_{j\in V}\exp(u_t\cdot v_j)$，
  只有 $j=k$ 这一项与 $v_k$ 有关，故
  \begin{equation}
  \frac{\partial g}{\partial v_k}
  =\frac{\partial\exp(u_t\cdot v_k)}{\partial v_k}
  =\exp(u_t\cdot v_k)\,u_t ,
  \end{equation}
  这里同样用了链式法则：外层 $\exp(\cdot)$ 在 $u_t\cdot v_k$ 处求值，
  内层 $u_t\cdot v_k$ 对 $v_k$ 的梯度是 $u_t$。
\end{itemize}
于是
\begin{equation}
\frac{\partial\log Z}{\partial v_k}
=f'(g)\cdot\frac{\partial g}{\partial v_k}
=\frac{1}{Z}\exp(u_t\cdot v_k)\,u_t
=P(k\mid t)\,u_t .
\end{equation}

\paragraph{(iii) 合并两种情况。}
把 (i)(ii) 两部分相加，并引入指示函数 $\mathbb{1}[k=c]$：
\begin{equation}
\boxed{\;
\frac{\partial y}{\partial v_k}
=\bigl(P(k\mid t)-\mathbb{1}[k=c]\bigr)\,u_t
=\begin{cases}
\bigl(P(k\mid t)-1\bigr)u_t, & k=c,\\[2pt]
P(k\mid t)\,u_t, & k\neq c .
\end{cases}\;}
\label{eq:grad-v}
\end{equation}

\textbf{对矩阵 $V$ 的梯度。} $\partial y/\partial V$ 的第 $k$ 行就是
$(\partial y/\partial v_k)^\top$，即
\begin{equation}
\left(\frac{\partial y}{\partial V}\right)_{k,:}
=\bigl(\hat y_k-(e_c)_k\bigr)u_t^\top,
\qquad\text{因此}\qquad
\boxed{\;\frac{\partial y}{\partial V}
=\bigl(\hat y-e_c\bigr)u_t^\top\;}
\label{eq:grad-V}
\end{equation}
其形状为 $|V|\times d$，与分母布局（梯度与变量同形）一致。

\subsection{题 1(c)：理解梯度（12 分）}

\paragraph{(i) “期望减观测”。}
在式 \eqref{eq:grad-u} 中，
$-v_c$ 是\textbf{观测}：它就是本训练对中真实出现的上下文词的向量，是数据的“正确答案”；
$\sum_{k\in V}P(k\mid t)v_k$ 是\textbf{期望}：它是模型当前预测分布 $P(\cdot\mid t)$ 下
上下文向量的期望，即 $\mathbb{E}_{k\sim P(\cdot\mid t)}[v_k]$，求期望的对象是模型预测的
上下文词分布。梯度等于“模型期望的上下文向量”减去“实际观测到的上下文向量”。

梯度为零当且仅当 $v_c=\sum_k P(k\mid t)v_k$，例如 $P(c\mid t)=1$（模型把概率全部押在
正确词上）。实际训练中\textbf{不可能}严格达到：softmax 输出对有限参数总是严格位于
$(0,1)$ 内，$P(c\mid t)<1$；而且同一个中心词在不同句子中会与不同上下文词配对，
这些约束互相冲突，不可能同时满足，因此梯度通常一直非零（呈噪声形式）。

\paragraph{(ii) 一步梯度下降的效果。}
固定 $u_t$，只用学习率 $\eta$ 更新 $v_c$。由式 \eqref{eq:grad-v}，
\begin{equation}
v_c^{\text{new}}
=v_c-\eta\bigl(P(c\mid t)-1\bigr)u_t
=v_c+\eta\bigl(1-P(c\mid t)\bigr)u_t .
\end{equation}
于是
\begin{equation}
u_t\cdot v_c^{\text{new}}
=u_t\cdot v_c+\eta\bigl(1-P(c\mid t)\bigr)\lVert u_t\rVert^2
\;\ge\; u_t\cdot v_c ,
\end{equation}
即正例对的点积（相似度）不会减小；只要 $u_t\neq 0$ 且 $P(c\mid t)<1$，它是严格增大的。
对 $k\neq c$，$v_k^{\text{new}}=v_k-\eta P(k\mid t)u_t$，于是
\begin{equation}
u_t\cdot v_k^{\text{new}}
=u_t\cdot v_k-\eta P(k\mid t)\lVert u_t\rVert^2 ,
\end{equation}
由于 softmax 下 $P(k\mid t)>0$，竞争词与中心词的点积严格减小。
也就是说，每次更新都把“中心词—真实上下文”拉近，同时把“中心词—被误认为可能的其它词”推远。
经常出现在相似上下文中的词，由于接受到的更新方向相近，会在向量空间中彼此靠近，
这正是“相似的词向量靠近”的来源。

\paragraph{(iii) 计算复杂度。}
对一个训练对，前向计算需要对词表中每个 $k$ 求点积 $u_t\cdot v_k$ 并做 softmax，
共 $O(|V|d)$ 次浮点乘法；由式 \eqref{eq:grad-u}、\eqref{eq:grad-v}，梯度对
$u_t$ 和所有 $v_k$ 的更新同样是 $O(|V|d)$。因此每个训练对的总代价为 $O(|V|d)$。

代入 $|V|=5\times10^5$、$d=300$：每个训练对约
$5\times10^5\times300=1.5\times10^8$ 次乘法；$10^8$ 个训练对共需约
\begin{equation}
1.5\times10^8\times10^8=1.5\times10^{16}
\end{equation}
次浮点乘法（把梯度与前向的常数因子算进去仍是 $10^{16}$ 量级）。

讲义第 37 页的负采样把“在 $|V|$ 个词上的多分类 softmax”换成
“1 个正例 + $K$ 个负例的二分类”：
\begin{equation}
y=-\log\sigma(u_t\cdot v_c)-\sum_{i=1}^{K}\mathbb{E}_{j\sim P(w)}
\log\sigma(-u_t\cdot v_j),
\end{equation}
每个训练对只需更新中心词与 $K+1$ 个上下文向量，代价降为 $O\bigl((K+1)d\bigr)$。
以 $K=5$ 为例，每个训练对只需 $6\times300=1800$ 次量级的乘法，
相比 $1.5\times10^8$ 约加速 $|V|/(K+1)\approx 8\times10^4$ 倍，
使在千万级到亿级词对的语料上训练成为可能。

% ============================================================
\section{第二部分 编程（55 分）}
% ============================================================

\subsection{环境与数据}

\textbf{机器配置}：Linux 服务器，Intel(R) Xeon(R) Silver 4210R CPU @ 2.40GHz，
40 核，62\,GB 内存，无 GPU 需求。Python 3.10（conda 环境 \texttt{nlp2026}），
依赖 \texttt{numpy}、\texttt{scipy}、\texttt{matplotlib}、\texttt{scikit-learn}、
\texttt{jieba}、\texttt{opencc-python-reimplemented}。

\textbf{语料}：中文维基百科过滤版 \texttt{pleisto/wikipedia-cn-20230720-filtered}
（约 25.5 万条词条，单个 JSON 约 500\,MB，正文在 \texttt{completion} 字段）。
\textbf{评测集}：\texttt{wordsim-240}、\texttt{wordsim-297}、\texttt{COS960}，
位于 \texttt{data/eval/}。

\textbf{训练路线}：路线 A，官方 C 代码
\href{https://github.com/tmikolov/word2vec}{tmikolov/word2vec}，commit \texttt{20c129a}，
由 \texttt{word2vec/train.sh} 自动 clone 到 \texttt{output/} 并编译，源码不进入仓库。

\subsection{题 2(a)：语料预处理（10 分）}

预处理脚本为 \texttt{word2vec/preprocess.py}，流程为：
读取 JSON（数组/JSONL 均可）取 \texttt{completion} 字段 $\to$ OpenCC 繁体转简体
$\to$ 小写并过滤标点、空白等无意义符号 $\to$ jieba 精确模式分词
$\to$ 每行一段、词间空格分隔写入 \texttt{data/corpus.txt}。
使用多进程（\texttt{multiprocessing.Pool}）按批量分词加速；分词后统计
$\text{min\_count}=5$ 的词表，并统计 3 个评测集在词表中的覆盖率与未登录词示例，
结果保存到 \texttt{word2vec/results/preprocess\_stats.json}，统计结果见表
\ref{tab:preprocess} 与表 \ref{tab:coverage}。复现命令：

\begin{verbatim}
conda activate nlp2026
python word2vec/preprocess.py
# 快速调试（只处理前 2000 条）：
python word2vec/preprocess.py --limit 2000 --output output/corpus_test.txt \
    --stats output/stats_test.json
\end{verbatim}

\begin{table}[htbp]
\centering
\caption{预处理统计（全量 25.5 万条词条）}
\label{tab:preprocess}
\begin{tabular}{lr}
\toprule
项目 & 数值 \\
\midrule
使用的词条数 & 254{,}547 \\
分词后总词数（token） & 82{,}854{,}433 \\
\texttt{min\_count}=5 的词表大小 & 414{,}976 \\
预处理耗时（秒） & 237.5（16 进程） \\
\bottomrule
\end{tabular}
\end{table}

表中词表大小按预处理脚本的统计口径给出；官方 word2vec 训练日志打印的
\texttt{Vocab size} 为 414{,}977，与计数器口径相差 1 个词，不影响训练设置。

\begin{table}[htbp]
\centering
\caption{评测集词覆盖率（未登录词比例）}
\label{tab:coverage}
\begin{tabular}{lrrr}
\toprule
数据集 & 唯一词数 & 覆盖词数 & 覆盖率 \\
\midrule
wordsim-240 & 327 & 326 & 99.69\% \\
wordsim-297 & 383 & 377 & 98.43\% \\
COS960 & 1{,}744 & 1{,}560 & 89.45\% \\
\bottomrule
\end{tabular}
\end{table}

\textbf{未登录词示例与分析}：以 \texttt{min\_count}=5 的词表为基准，wordsim-240 仅
1 个唯一词未覆盖，wordsim-297 有 6 个，COS960 有 184 个。对典型未登录词回查全量语料
\texttt{data/corpus.txt} 的词频后，可以把未登录词分成四类：

\begin{itemize}
  \item \textbf{被 jieba 切碎}：如“日本人”“美国人”“社会调查”“相似度”在语料中并不以
  整词出现，而是被切成“日本/人”“美国/人”“社会/调查”“相似/度”等高频片段；
  \item \textbf{低频词被 \texttt{min\_count} 过滤}：如“前情”（4 次）、“做秀”（1 次）、
  “倔犟”（1 次）、“口涎”（2 次）、“奸笑”（4 次）、“吸顶灯”（1 次），出现次数不足 5 次；
  \item \textbf{语料中确实没有}：如“乖觉”“反剪”“前轮胎”“后轮胎”“心签名”等在全量语料中
  出现 0 次；
  \item \textbf{英文大小写不一致}：wordsim-297 的 \texttt{OPEC} 在预处理中被统一小写为
  \texttt{opec}（语料中出现 28 次），但评测脚本按原样精确匹配，因此被误判为未登录
  （评测时对词对统一小写即可覆盖）。
\end{itemize}

\textbf{能否用 \texttt{jieba.load\_userdict} 提高覆盖率}：可以，把评测集词表加入用户词典后，
jieba 会优先把这些词切成整体，从而减少“被切碎”的 OOV。但这样做有两个问题：
(i) 评测集本身被用于分词，等于把测试信息泄漏进训练流程，评出的覆盖率与相似度会偏乐观，
不再反映真实泛化；(ii) 用户词典只解决“分词粒度不一致”造成的假 OOV，
如果语料中确实没有该词，仍然无法覆盖。因此正式实验中不建议用评测集构建用户词典，
可在报告里作为对照实验说明。

\subsection{题 2(b)：训练（12 分）}

本实验采用路线 A，训练 skip-gram + 负采样词向量
\citep{mikolov2013efficient,mikolov2013distributed}。
基线超参数：维度 $d=200$、窗口 $m=5$、负例数 $K=5$、下采样阈值 $10^{-4}$、
\texttt{min\_count}=5、轮数 5，线程数固定为 8。训练命令（由 \texttt{train.sh} 执行）：

\begin{verbatim}
output/word2vec_official/word2vec \
  -train data/corpus.txt -output output/sg_w5.txt \
  -cbow 0 -size 200 -window 5 -negative 5 -hs 0 \
  -sample 1e-4 -min-count 5 -iter 5 -threads 8 -binary 0
\end{verbatim}

官方 C 代码不含随机种子参数，训练过程依赖负采样表等内部随机性；
本实验固定线程数与迭代轮数以保证设置可比。机器配置与完整命令如上，
基线（$m=5$）训练耗时 1648 秒（约 27.5 分钟，含加载与写出时间），
窗口对比各设置的耗时见表 \ref{tab:hparams}。

训练日志保存在 \texttt{output/logs/}。官方 C 代码不输出逐轮损失，但会按训练进度记录
当前学习率（\texttt{Alpha} 与 \texttt{Progress}），\texttt{visualize.py --which lr}
据此绘制学习率随训练进度的线性衰减曲线，见图 \ref{fig:lr}。

\figfile{0.8\linewidth}{lr_curve.png}{基线训练的学习率随训练进度的线性衰减}{fig:lr}

\subsection{题 2(c)：内在评测（15 分）}

评测脚本 \texttt{word2vec/evaluate.py} 自行实现了余弦相似度（对词向量矩阵做
$\ell_2$ 归一化后矩阵乘）、最近邻检索（一次矩阵--向量乘法，排除查询词自身）
与 Spearman 相关（调用 \texttt{scipy.stats.spearmanr}），
未使用 gensim 的评测函数。三个评测集分别来自 \citet{jin2012semeval}
（wordsim-240/297）与 \citet{huang2019cos960}（COS960）。运行：

\begin{verbatim}
python word2vec/evaluate.py --vectors output/sg_w5.txt
\end{verbatim}

\paragraph{(i) 词相似度。}
结果见表 \ref{tab:similarity}。

\begin{table}[htbp]
\centering
\caption{基线词向量的词相似度结果（OOV 词对跳过）}
\label{tab:similarity}
\begin{tabular}{lrrr}
\toprule
数据集 & 总对数 & 覆盖对数 & Spearman $\rho$ \\
\midrule
wordsim-240 & 240 & 239 & 0.5457 \\
wordsim-297 & 297 & 292 & 0.5961 \\
COS960 & 960 & 801 & 0.4453 \\
\bottomrule
\end{tabular}
\end{table}

三个数据集上均取得中等偏上的正相关（$\rho\approx0.45$--$0.60$）。wordsim 系列衡量
“相关性”（如“李白—诗”），而 COS960 只衡量“近义程度”，且包含大量仅相关而非近义
的词对，因此 $\rho$ 略低是符合预期的。

\paragraph{(ii) 最近邻。}
固定查询词：北京、语言、苹果、计算机、足球、孔子、春节、快乐；
自选查询词：小米、银行、同志、打（多义词）。各查询词余弦最近的 10 个词如下
（括号内为余弦值，完整结果见 \texttt{results/neighbors\_sg\_w5.tsv}）：

\begin{itemize}
  \item \textbf{北京}：北京市 0.698、上海 0.678、沈阳 0.668、北京城 0.652、
  天津 0.633、燕翔 0.629、呼和浩特 0.628、工人体育场 0.624、昆仑饭店 0.620、海淀区 0.616
  \item \textbf{语言}：语法 0.768、书面语言 0.759、语等 0.752、第二语言 0.741、
  母语 0.740、语是 0.739、语系 0.738、语法结构 0.738、语族 0.737、语 0.734
  \item \textbf{苹果}：黑莓 0.649、坚果 0.639、小米 0.632、火龙果 0.624、
  脆片 0.621、南瓜 0.618、主打产品 0.610、草莓 0.604、洋葱 0.603、鳄梨 0.598
  \item \textbf{计算机}：图形学 0.799、计算机网络 0.787、人机交互 0.785、
  计算机系统 0.785、信号处理 0.773、应用软件 0.772、计算机辅助 0.769、
  分布式计算 0.761、电脑 0.759、电子计算机 0.758
  \item \textbf{足球}：甲级联赛 0.777、俱乐部 0.735、篮球 0.733、足球联赛 0.732、
  国际足球 0.731、联赛 0.725、球队 0.717、司职 0.714、五人制 0.709、青年队 0.705
  \item \textbf{孔子}：颜回 0.654、孟子 0.626、曾子 0.615、孔丘 0.613、子思 0.613、
  十哲 0.603、confuciusornis 0.601、圣贤 0.599、颜何 0.598、家语 0.595
  \item \textbf{春节}：联欢晚会 0.789、元宵 0.738、央视春晚 0.718、大年初一 0.683、
  连休 0.677、春节假期 0.673、除夕夜 0.672、郭冬临 0.671、中秋节 0.665、春节晚会 0.660
  \item \textbf{快乐}：锋尚 0.663、男声 0.659、幸福 0.648、爸爸妈妈 0.647、敢不敢 0.643、
  强合辑 0.640、好时光 0.637、谁动 0.634、相亲相爱 0.634、一生一世 0.626
  \item \textbf{小米}：红米 0.719、miui 0.692、pro2s 0.646、smartisan 0.642、苹果 0.632、
  新品 0.625、oneplus 0.623、手机 0.607、雷军 0.606、旗舰机 0.606
  \item \textbf{银行}：商业银行 0.815、分行 0.813、总行 0.766、金融机构 0.747、
  兴业银行 0.743、中国银行 0.739、储蓄银行 0.734、bank 0.733、该行 0.732、国家银行 0.726
  \item \textbf{同志}：党外 0.581、祖国统一 0.574、老战士 0.562、平权 0.559、
  恳谈会 0.552、积极分子 0.550、少先队 0.547、女同性恋 0.546、台湾同胞 0.542、保卫和平 0.542
  \item \textbf{打}：席面 0.651、平飞安 0.650、打打 0.643、一个打 0.633、廖健富 0.632、
  安打及 0.632、点圈 0.626、打以 0.626、lofton 0.624、球击 0.623
\end{itemize}

\textbf{“意外”近邻与解释}：

(i) \textbf{苹果}的近邻以水果为主（黑莓、火龙果、草莓、鳄梨），却混入了“小米”；
而查询\textbf{小米}的近邻几乎全是手机品牌（红米、miui、雷军、苹果、oneplus）。
“苹果”“小米”“黑莓”都是“水果/科技品牌”双义词：词向量把两个义项的语境平均在同一个
向量里，而维基语料中“小米公司”的语料远多于“小米（谷物）”，因此小米向量偏向科技义项，
苹果向量则偏水果义项，但两者仍因共同义项互相靠近。

(ii) \textbf{孔子}的近邻中出现拉丁学名 \texttt{confuciusornis}（孔子鸟），
它在语料中仅出现 9 次，几乎只与“孔子”共现（该属名用以纪念孔子），
这是“共现联想”而非语义相似造成的近邻。

(iii) \textbf{春节}的近邻包含“郭冬临”（仅出现 18 次的演员名），
原因是他是央视春晚的常客，二者上下文高度重叠；\textbf{快乐}的近邻
（锋尚、强合辑、谁动、相亲相爱等）多来自“快乐男声”“快乐大本营”等节目与歌曲名，
属于固定搭配造成的伪近邻。

(iv) \textbf{打}作为高频多义动词，近邻大量是分词碎片（“安打及”“打以”“一个打”）
和棒球术语（安打、球击、lofton、廖健富），说明高频虚化动词的向量被多种语境稀释，
单独一个向量难以承载其全部义项。相比之下，语义较单一的“银行”“计算机”
近邻质量很高（商业银行、分行、总行；图形学、人机交互、电脑）。

\paragraph{(iii) 误差分析。}
按模型余弦排名与人工打分排名的名次差（\texttt{rankdata} 平均秩）排序，
COS960 上差异最大的 5 对（保存于 \texttt{results/cos960\_errors\_sg\_w5.tsv}）及原因分析：

\begin{itemize}
  \item \textbf{聪明绝顶—胆小怕事}（人工 0.00，余弦 0.844，模型第 47 名 / 人工第 746 名）：
  两者都是描写人物性格的四字成语，常出现在“他是一个……的人”等相同句式中，
  语料中仅 10 与 16 次。模型只看到语境相似，而 COS960 标注的是“近义”，
  二者只是“相关”，属于数据集语义与共现语义的错位。
  \item \textbf{困倦—烦燥}（人工 0.00，余弦 0.843）：都是身体/情绪状态词，
  共享“感到~”“有些~”语境；另外“烦燥”是“烦躁”的异形写法（语料中仅 7 次，
  而“烦躁”有 73 次），低频异形词的向量估计本就不稳。
  \item \textbf{球队—赛季}（人工 0.00，余弦 0.807）：体育领域高相关词
  （语料中分别出现 36\,085 与 20\,261 次），上下文高度重叠，但“相关”不等于“近义”，
  再次体现 COS960 与 wordsim 的差别。
  \item \textbf{悲鸣—哀鸣}（人工 3.87，余弦 0.462，模型第 678 名 / 人工第 41 名）：
  真正近义的书面语词，但“哀鸣”在语料中仅出现 17 次（“悲鸣”128 次），
  低频词的向量方向不稳定，且两词常在文体的不同上下文中单独出现。
  \item \textbf{父母—爸妈}（人工 4.00，余弦 0.514，模型第 617 名 / 人工第 11 名）：
  语义近义但语体不同，“爸妈”口语化且仅出现 65 次（“父母”6931 次），
  二者上下文分布有明显差异（“父母”多用于书面与正式叙述）。
\end{itemize}

总体来看，模型错误的两个主要来源是：(1) 训练目标捕获的是“共现/相关”而非
“近义”，因此“球队—赛季”“聪明绝顶—胆小怕事”这类相关词对被判为高度相似；
(2) 低频词与异形词（哀鸣、爸妈、烦燥）的向量估计不稳定。

\subsection{题 2(d)：超参数对比（10 分）}

选择窗口大小对比：$m\in\{2,5,10\}$（$m=5$ 为基线），其余超参数保持不变，
即 \texttt{bash word2vec/train.sh window}。各设置在 3 个评测集上的结果与训练耗时
见表 \ref{tab:hparams} 与图 \ref{fig:hparams}：

\begin{table}[htbp]
\centering
\caption{不同窗口大小的结果与训练耗时（秒）}
\label{tab:hparams}
\begin{tabular}{lrrrr}
\toprule
设置 & $\rho$ (240) & $\rho$ (297) & $\rho$ (COS960) & 训练耗时 \\
\midrule
$m=2$ & 0.5134 & 0.5668 & 0.4587 & $\approx$900（估算） \\
$m=5$（基线） & 0.5457 & 0.5961 & 0.4453 & 1648 \\
$m=10$ & 0.5484 & 0.5855 & 0.4158 & 2480 \\
\bottomrule
\end{tabular}
\end{table}

其中 $m=2$ 的耗时未由 \texttt{train.sh} 精确记录（当时脚本尚未加入计时），
按训练日志的吞吐量（约 $6.2\times10^4$ 词/线程/秒）估算；其余为脚本记录值。

\figfile{0.8\linewidth}{hparams.png}{窗口大小对 3 个评测集 Spearman $\rho$ 的影响}{fig:hparams}

分析：窗口增大时，“相关性”与“近义性”的变化方向并不一致：
wordsim-240 随窗口增大单调上升（$0.5134\to0.5457\to0.5484$），
wordsim-297 在 $m=5$ 达到峰值（$0.5668\to0.5961\to0.5855$），
而 COS960 随窗口增大单调下降（$0.4587\to0.4453\to0.4158$）。
这与 \citet{levy2015improving} 的结论一致：大窗口使向量编码更多“话题/领域相关”，
有利于 wordsim 这类相关性评测；小窗口聚焦紧邻语境，更能捕获可相互替换的
“同义/同功能”关系，因而更有利于 COS960 的近义评测。训练耗时随窗口近似线性增长，
$m=10$ 为 2480 秒，约为 $m=5$（1648 秒）的 1.5 倍。

不同设置下近邻词的变化以“足球”最为典型：$m=2$ 时前 5 个近邻为
篮球 0.74、排球 0.70、足球运动 0.67、五人制 0.67、足球队 0.66，
以“同类可替换的球类运动”为主；$m=10$ 时为
甲级联赛 0.83、足球联赛 0.80、俱乐部 0.79、联赛 0.77、球队 0.77，
全部是同一话题领域的联赛/俱乐部词，余弦也更高。
这说明窗口越大，模型越偏“话题相关”，越小越偏“功能/类别相似”。
“苹果”也类似：$m=2$ 的近邻是橘子、马铃薯、柑橘、椰树（同类食物），
$m=10$ 时为黑莓、水果、南瓜、软果（仍偏食物，但更偏共现话题）。

\subsection{题 2(e)：可视化（8 分）}

可视化脚本 \texttt{word2vec/visualize.py} 选取 7 个语义类别
（地区、动物、颜色、朝代、职业、体育、学科）共 81 个词，
分别用 PCA 与 t-SNE 降到 2 维并按类别着色，结果见图 \ref{fig:pca} 与图 \ref{fig:tsne}。
这 81 个词在 \texttt{data/corpus.txt} 中的词频均不低于 5，
实际全部保留在词表中（81/81），满足 $\ge 5$ 类、$\ge 50$ 词的要求。

\figfile{0.85\linewidth}{pca.png}{语义类别的 PCA 二维散点图}{fig:pca}
\figfile{0.85\linewidth}{tsne.png}{语义类别的 t-SNE 二维散点图}{fig:tsne}
\figfile{0.9\linewidth}{neighbors.png}{查询词（星号）及其 top-10 近邻的 t-SNE 图}{fig:neighbors}

\textbf{讨论}：

(i) \textbf{PCA 与 t-SNE 的类别分离度}：按类别计算二维坐标的轮廓系数，
PCA 为 0.425，t-SNE 为 0.644（perplexity $=27$、seed $=42$），t-SNE 明显分得更开。
图 \ref{fig:pca} 中颜色、地区等类别已经分开，但体育、职业、学科三类在下方互相
交叠；图 \ref{fig:tsne} 中 7 个类别几乎各自成簇，只有职业（教师、记者、律师等）
与学科略有交叠，符合“教师/记者/律师”与学科语境相关的直觉。

(ii) \textbf{t-SNE 簇间距离不可直接解读为语义距离}：t-SNE 只保留局部邻域结构，
簇间相对位置与距离随参数变化很大。对比图 \ref{fig:tsne}（perplexity $=27$）
与图 \ref{fig:tsne-p10}（perplexity $=10$）：两次运行中“地区”簇分别位于画布
左侧与右侧，“颜色”簇分别在上方与右下，但各簇内部始终紧密；因此只能解读
“簇内/近邻关系”，不能把两个簇的间距当作语义相似度。

(iii) \textbf{稳定性}：改变随机种子（seed $42\to0$，图 \ref{fig:tsne-s0}）后，
全局布局再次整体旋转/平移，轮廓系数从 0.644 变为 0.631；改变 perplexity 后为
0.684。类别内部始终聚在一起，但个别词的位置（如“北京”在地区簇与朝代簇之间的
位置）会在不同运行间移动，说明 t-SNE 用于观察局部结构是可靠的，
但不宜过度解读具体坐标。

(iv) \textbf{近邻群是否重叠}：图 \ref{fig:neighbors} 中，北京、足球、春节、
计算机、孔子、苹果 6 组基本互不重叠，各自内部紧凑：北京组包含“天津、上海、
工人体育场、昆仑饭店”等城市与地标；孔子组（颜回、孟子、曾子）独立于其他组。
唯一的“跨界”是苹果组混入“小米”，与题 2(c)(ii) 中观察到的
“水果/科技品牌”双义现象一致。

\figfile{0.85\linewidth}{tsne_p10.png}{perplexity $=10$ 时的 t-SNE 图（与图 \ref{fig:tsne} 对比簇间位置变化）}{fig:tsne-p10}
\figfile{0.85\linewidth}{tsne_s0.png}{随机种子 $=0$ 时的 t-SNE 图（与图 \ref{fig:tsne} 对比稳定性）}{fig:tsne-s0}

% ============================================================
\section{第三部分 提交信息}
% ============================================================

\begin{itemize}
  \item 仓库地址：\todo{https://github.com/.../nlp-fall2026-lab4}
  \item 最后一次提交的 commit id：\todo{ }
  \item 训练路线：A（tmikolov/word2vec @ 20c129a）
\end{itemize}

\bibliographystyle{plainnat}
\bibliography{refs}

\end{document}
